\documentclass[10pt,a4paper]{amsart}
\usepackage[margin=19mm]{geometry}
\usepackage{amsmath,amssymb,mathtools}
\usepackage[scr=rsfs,cal=euler]{mathalfa}
\usepackage{xfrac}
\usepackage[unicode,hidelinks,bookmarks=false]{hyperref}
\newcommand{\ie}{i.e.\ }
\newcommand{\cf}{cf.\ }
\newcommand{\resp}{respectively\ }
\newcommand{\Subsection}{\subsection}
\providecommand{\nobreakdash}{\nobreak\hbox{-}}
\setlength{\emergencystretch}{2em}
\multlinegap0pt
\usepackage[shortlabels]{enumitem}
\usepackage{amssymb}
\usepackage{xcolor}
\newtheorem{definition}{Definition}
\newtheorem{lemma}{Lemma}
\newtheorem{corollary}{Corollary}
\newtheorem{proposition}{Proposition}
\def\N{{\mathbb{N}}}
\def\R{{\mathbb{R}}}
\newcommand{\RRP}{\mathsf{RRP}}
\newcommand{\cD}{\mathcal{D}}
\newcommand{\e}{\varepsilon}
\renewcommand{\phi}{\varphi}
\title{Full measure universality for Cantor Sets}
\author{Pablo Shmerkin, Alexia Yavicoli}
\date{Source: arXiv:2503.21079v2 (CC BY 4.0)}
\begin{document}
\maketitle

\noindent\emph{Source and licence.} Full measure universality for Cantor Sets. Version arXiv:2503.21079v2 (CC BY 4.0), distributed by arXiv under the Creative Commons Attribution 4.0 International licence, http://creativecommons.org/licenses/by/4.0/, as recorded in the arXiv licence metadata for this exact version and shown on the abstract page https://arxiv.org/abs/2503.21079v2. Full licence notice: \texttt{licenses/arxiv-2503.21079-CC-BY-4.0.txt}.\par\smallskip

\noindent\textit{Editorial scope.} Counted unit: the article's proposition proposition (source0.tex lines 924--930). The article's complete original proof of it is delivered with it (source0.tex lines 931--992, one \texttt{proof} environment of 62 lines). No mathematics is written, repaired or re-derived: the statement and the proof are the article's own lines, in the article's own notation, and no retained line is substituted or re-flowed.  Retained uncounted as the source-local context the proof needs: lines 270--286; lines 349--358; lines 413--418; lines 864--875.  Nothing was omitted from any retained span, so the package carries no omission record.  No retained line contains a citation, so the wrapper carries no bibliography and no bibliography entry.  No retained line contains a figure environment, an image-inclusion command or drawing code, so no image is delivered and the package declares no asset dependency.  Editorial formatting only, in addition: an AMS \texttt{amsart} preamble that restates the article's own declarations and macros used by the retained lines, namely the environments \texttt{definition}, \texttt{lemma}, \texttt{corollary}, \texttt{proposition} on separate counters, together with the standard packages those lines need.  The retained environments print in this document as Definition 1, Lemma 1, Corollary 1, Proposition 1 in order of appearance, whereas the article prints the same items with the numbering of its own full document.  The complete native source of the article as distributed in the arXiv e-print for this exact version is bundled unchanged as \texttt{sources/175462/source0.tex}.
\begin{definition} \label{def:large-set}
    Let $A\subset\R^d$ be a Borel set and let $N(\delta):(0,1]\to\N$ be a function such that
    \[
        \lim_{\delta\to 0^{+}} N(\delta) = \infty.
    \]

    We say that $A$ is \emph{$N$-large} if there exists a compact set $A'\subset A$ such that for every cube $Q$ such that $A'\cap Q^{\circ}\neq\emptyset$,
    \[
       \limsup_{\delta\to 0^{+}} \frac{|A'\cap Q^{\circ}|_{\delta}}{N(\delta)} \ge 1.
    \]

    If $\delta_k, N_k$ are sequences in $(0,1]$ and $\N$ converging to $0$ and $\infty$ respectively, we say that $A$ is \emph{$(N_k,\delta_k)$-uniformly large} if it contains a compact set $A'$ with the following property:
    \[
        |A'\cap Q|_{\delta_k} \ge N_k
    \]
    for all cubes $Q\in\mathcal{D}_k(A')$.
\end{definition}

\begin{lemma} \label{lem:Hausdorff-unif-large}
    Let $\phi\in\Phi$, let $\eta>0$, and let $\delta_k$ be a sequence such that
    \[
        \lim_{k\to\infty} 2^{3kd+1} \phi(\delta_k) = 0.
    \]
    Then, any set $A\subset [0,1]^d$ with $\mathcal{H}^\phi(A)\ge \eta$ is $(N_k,\delta_k)$-uniformly large for
    \[
      N_k= \frac{\eta}{2^{3kd+1} \cdot \phi(\delta_k)}.
    \]
\end{lemma}

\begin{lemma} \label{lem:nonemptyinterior}
    Suppose $(A,\mathcal{F})$ has the $\RRP$ property with initial rectangle $R$. Then for every $a_0\in A$ there is a compact set $K$ with $|K|=0$ such that
    \begin{equation} \label{eq:nonemptyinterior}
        f(a_0)+R\subset f(A)+K \quad\text{for all } f\in\mathcal{F}.
    \end{equation}
\end{lemma}

\begin{corollary} \label{cor:large-sum-R^d}
    For every $\eta,\e\in (0,1/3]$, there is $N=N(\eta,\e)\in\N$ such the following holds for all small enough $\delta>0$.

    Fix a dyadic scale $2^{-m'}$ and a cube $Q\in\cD_{m'}$. There exists a set $B\subset 4\cdot Q$ with $|B|\le \eta\cdot |Q|$, such that
    \[
        \bigl|(a_0+Q)\cap (A+B)\bigr| \ge (1-\e)|Q|
    \]
    for all  $A\subset [0,2^{-m'}]^d$ with  $|A|_{\delta}\ge N$ and all $a_0\in A$.


    Moreover, $B$ can be taken to be a union of closed dyadic cubes of the same size, depending only on $d,\eta,\e,\delta,m'$.
\end{corollary}

\begin{proposition}
    Let $\phi\in\Phi$ and $\eta,\e>0$. There exists a compact set $B$ (depending only on $\phi$, $\eta$, $\e$ and the ambient dimension $d$) with $|B|=0$, such that
    \[
        \bigl|[0,1]^d\cap (A+B)\bigr| \ge 1-\e,
    \]
    for all Borel sets $A\subset [0,1]^d$ with $\mathcal{H}^{\phi}(A)>\eta$.
\end{proposition}

\begin{proof}
    The proof is similar to that of Lemma \ref{lem:nonemptyinterior}. Invoking Lemma \ref{lem:Hausdorff-unif-large}, we may assume that $A$ is $(N_k,\delta_k)$-uniformly large for some $N_k\to\infty$ and $\delta_k$ that depend only on $\phi$ and $\eta$.

    Let $\mathcal{A}$ denote the family of all sets $A$ with the following property: for any $k$ and any $Q\in\mathcal{D}_k(A)$, we have
    \begin{equation} \label{eq:good-assumption}
        |A\cap Q|_{\delta_k} \ge N_k.
    \end{equation}
    By the definition of uniformly large (Definition \ref{def:large-set}), it is enough to prove the claim for $A\in\mathcal{A}$.

    We will inductively construct a sequence of sets $B_j\subset [-1,1]^d$ and a strictly increasing sequence of integers $m_j\in\N$ such that:
    \begin{enumerate}[label=(\alph*)]
        \item \label{it:a-PM} $B_j$ is a union of closed cubes $\{Q^{(j)}_s\}_{s\in S_j}$, where $Q^{(j)}_s\in\cD_{m_j}$.
        \item \label{it:b-PM} $B_{j+1}\subset B'_j$, where
        \[
            B'_j = \bigcup_{s\in S_j} 4 Q ^{(j)}_s.
        \]
        \item \label{it:c-PM} $|B_j|\le  2^{-j}$,
        \item \label{it:d-PM} For all $A\in\mathcal{A}$, there is a finite subset $A'_j\subset A$ such that $\left| [0,1]^d\setminus (A'_j+B_j)\right|\le (1-2^{-j})\e$.

    \end{enumerate}
    Assuming the construction can be carried out, let
    \[
        B''_j = \bigcup_{s\in S_j} 6 Q^{(j)}_s,
    \]
    and note from \ref{it:a-PM}--\ref{it:b-PM} that $B''_j$ is a nested sequence. Let $B=\bigcap_{j=1}^\infty B''_j$. Then $|B|=0$ by \ref{it:c-PM}, since $|B''_j|\le 6^d |B_j|$. If $A\in\mathcal{A}$, we have
    \[
        \bigl|[0,1]^d\setminus (A+B)\bigr| = \lim_{j\to\infty} \left|[0,1]^d\setminus (A+B''_j)\right| \le \e.
    \]
    Thus, the existence of the claimed set $B$ follows from the construction.

    Set $B_0=[-1,1]^d$, $m_0=0$ and $A'_0$ given by an arbitrary element of $A$ for all $A\in\mathcal{A}$. Suppose $B_j$, $m_j$ have been constructed. Suppressing the dependence on $j$ for simplicity, let
    \[
      B_j = \bigcup_{s\in S} Q_s, \quad Q_s\in\mathcal{D}_{m}.
    \]
    Let
    \[
        N' = N(1/2,2^{-(j+1)}\e),
    \]
    where $N$ is the function from Corollary \ref{cor:large-sum-R^d}.

    Since $N_k\to\infty$ and $\delta_k\to 0$, we can pick $k\ge m$ so that $N_k\ge N'$ and $\delta_k$ is small enough that the conclusion of Corollary \ref{cor:large-sum-R^d} holds.

    Let
    \[
        B_{j+1} = \bigcup_{s\in S} B'_s,
    \]
    where $B'_s$ is the outcome of applying Corollary \ref{cor:large-sum-R^d} to $Q_s$.  Claims \ref{it:a-PM}, \ref{it:b-PM} and \ref{it:c-PM} hold by Corollary \ref{cor:large-sum-R^d} and the choice of $N'$.

    Fix $a_0\in A'$. Then, since $A\in\mathcal{A}$ and $k\ge m$, we know that $|A\cap R_0|_{\delta}\ge N'$, where $R_0$ is the dyadic cube in $\cD_{m}$ containing $a_0$. Let $A''(a_0)$ be a finite subset of $A\cap R_0$ such that $|A''(a_0)|_{\delta}\ge N'$. Then, by Corollary \ref{cor:large-sum-R^d},
    \[
        \bigl| (a_0+Q_s)\cap (A''(a_0) + B'_s) \bigr| \ge \left(1-2^{-j-1}\e\right)|Q_s|.
    \]
    Adding up over all $s$, we get that
    \[
        \bigl| (A'_j  + B_j) \setminus (A'_{j+1} + B_{j+1})\bigr| \le 2^{-(j+1)}\e.
    \]
    We conclude from the inductive assumption \ref{it:d-PM} that
    \[
        \bigl| [0,1]^d\setminus (A'_{j+1}+B_{j+1})\bigr| \le (1-2^{-j})\e + 2^{-(j+1)}\e = (1-2^{-(j+1)})\e.
    \]
    This shows that \ref{it:d-PM} continues to hold for $j+1$, and the induction is complete.
\end{proof}

\end{document}
